\SourceLeaf{58}
\noindent Here too the way to the general case is already open. If, as we may, we equate the expressions multiplying $\sin(h+1)\varphi$ on both sides, we obtain

\[
\frac{\nu nA^{(h)}}{2}-\frac{\nu nA^{(h+2)}}{2}=(h+1)A^{(h+1)}+\frac{nhA^{(h)}}{2}+\frac{n(h+2)A^{(h+2)}}{2}.
\]

\noindent Solving the equation gives

\[
A^{(h+2)}=\frac{n(\nu-h)A^{(h)}-2(h+1)A^{(h+1)}}{n(\nu+h+2)},
\]

\noindent This general formula contains the law for every coefficient in the series we have posited, with the exception of what would result for $h=0$, namely $A^{(2)}$, which has already been found by a separate investigation. It gives the coefficient from the preceding two by a simple calculation. For example,

\[
A^{(2)}=\frac{2(\nu nA-A^{(1)})}{n(\nu+2)};\qquad A^{(3)}=\frac{n(\nu-1)A^{(1)}-4A^{(2)}}{n(\nu+3)};
\]

\[
A^{(4)}=\frac{n(\nu-2)A^{(2)}-6A^{(3)}}{n(\nu+4)};\quad\mbox{etc.}
\]

$\delta$) However, when $\nu$ is a negative number, a denominator can vanish and leave a coefficient undefined. In that case we must return to the infinite series given earlier. We must do the same for the first two coefficients $A$ and $A^{(1)}$, which the recurrence scheme already assumes to be known. By a particular device, however, $A^{(1)}$ can be expressed in terms of $A$. For we had

\[
A=1+.\,.\,.\,.+\frac{1}{2^{2h}}{}^{\nu}\mathfrak{B}^{(2h)}.{}^{2h}\mathfrak{B}^{(h)}n^{2h}+.\,.\,.,
\]

\[
A^{(1)}={}^{\nu}\mathfrak{B}^{(1)}n+.\,.\,.+\frac{1}{2^{2h}}{}^{\nu}\mathfrak{B}^{(2h+1)}.{}^{2h+1}\mathfrak{B}^{(h+1)}n^{2h+1}+.\,.\,..
\]

\noindent If we take $2nA+A^{(1)}$, the initial term of this series is

\[
=({}^{\nu}\mathfrak{B}^{(1)}+2)n=(\nu+2)n,
\]

\noindent and the general term $h^{\text{tus}}$ is

\[
2n . \frac{1}{2^{2h}}{}^{\nu}\mathfrak{B}^{(2h)}.{}^{2h}\mathfrak{B}^{(h)}n^{2h}+\frac{1}{2^{2h}}{}^{\nu}\mathfrak{B}^{(2h+1)}.{}^{2h+1}\mathfrak{B}^{(h+1)}n^{2h+1},
\]

\noindent Since

\[
{}^{\nu}\mathfrak{B}^{(2h+1)}={}^{\nu}\mathfrak{B}^{(2h)} . \frac{\nu-2h}{2h+1}\qquad\mbox{and}\qquad {}^{2h+1}\mathfrak{B}^{(h+1)}={}^{2h}\mathfrak{B}^{(h)} . \frac{2h+1}{h+1},
\]

\noindent this term simplifies to

\[
\frac{1}{2^{2h}}{}^{\nu}\mathfrak{B}^{(2h)}.{}^{2h}\mathfrak{B}^{(h)}\left(2+\frac{\nu-2h}{2h+1} . \frac{2h+1}{h+1}\right)n^{2h+1}
\]

